
\lhead{\includegraphics[height=1cm]{Logos And Group/LHCb_Logo.png}}
\chead{Hinweise zur Benutzung des Quark-Puzzles}
\rhead{\includegraphics[height=1cm]{Logos And Group/University of Bonn_logo.png}}
\begin{document}
\hfill\textit{Zur Verwendung für Vermittler*innen und Lehrpersonen}
\section*{Hinweise zum Quark-Puzzle (erste Version)}

Das Quark-Puzzle ist als Lehrmittel der Hadronenphysik konzipiert und soll Schülerinnen, Schülern und Interessierten das eigenständige Entdecken von Eigenschaften der Starken Wechselwirkung und Hadronen ermöglichen. Es bildet außerdem eine Anwendungsmöglichkeit zum Erlernen der Quarks des Standardmodells der Teilchenphysik. 

Das Puzzle besteht aus insgesamt 36 unterschiedlichen Teilen (Quarks), in den Ausführungen (Flavours) up, down, strange, charm, top, bottom und den entsprechenden Antiteilchen. Zur Veranschaulichung der Ladung der Starken Wechselwirkungen (Farbladung) liegen alle Quarks in drei Versionen vor, die falls möglich, in den jeweiligen Farben/Antifarben (rot, blau, grün für Materie; Antirot (Cyan), Antigrün (Magenta) und Antiblau (Gelb) für Antimaterie) ausgedruckt werden sollten.  
	Der Steckmechanismus ist so ausgewählt, dass nur farbneutrale Hadronen gebaut werden können, dass bedeutet, folgende Kombinationen sind möglich:
\begin{table}[h]
    \centering
       \caption[Übersicht über die möglichen Farbkombinationen des Quark-Puzzles  und sich daraus ergebene Modelle farbneutraler Hadronen.]{Übersicht über die möglichen Farbkombinationen des Quark-Puzzles und sich daraus ergebene Modelle farbneutraler Hadronen.}
    \label{tab:QuarkPuzzle}
    \begin{tabular}{ccccc|c}
         & & Kombinationen &&   & Hadron\\ \hline
        Blau &-& Rot &-& Grün  & Baryon\\
        Antiblau &-& Antirot &-& Antigrün & Antibaryon\\
        Blau &-& Antiblau &&& Meson\\
        Rot &-& Antirot&&& Meson \\
        Grün &-& Antigrün&&& Meson
    \end{tabular}
    \end{table}

Auf einem einzelnen Quark sind folgende Informationen verzeichnet (s. auch Bilder):

  \begin{description}
    \item[Generationen]werden auf einer Seite über Formelemente repräsentiert, dabei steht die Form des Zylinders für die erste, die des Tetraeders für die zweite und die des Quaders für die dritte Generation.
    \item[Farbneutralität:]Es lassen sich nur farbneutrale Hadronen formen, der Flavour besitzt keinen Einfluss auf den Steckmechanismus.
    \item[\n{Der} Flavour]ist auf einer weiteren Seite eingraviert, wobei Antiquarks einen Strich über dem Buchstaben besitzen
    \item[\n{Die} elektrische Ladung]der Quarks ist auf einer weiteren Seite als Wert eingraviert.
    \item[Farbladungen]werden anhand der Druckfarbe repräsentiert.
    \item[Top-Quarks]besitzen keine Verbindungen, sodass sie nicht zu einem Hadron zusammengebaut werden können. Hierdurch wird veranschaulicht, dass Hadronen mit $t$-Flavour nicht beobachtet werden.
\end{description}

\subsection*{Einbettung in unterrichtliche Kontexte/Masterclasses}

Das Quark-Puzzle verfolgt neben den nachfolgenden Erklärungen zu fachlichen Aspekten vor allem Ziele der Veranschaulichung. Ein Ziel ist, dass ein abstrakt wirkendes Thema, mit zahlreichen unbekannten Vokabeln, spielerisch und infolge mit Interesse erkundet wird. Es soll so ein grundlegender Eindruck/mentale Repräsentation über die physikalischen Prinzipien und Begriffe entstehen, auf die infolge inhaltlich systematisch aufgebaut werden kann. 

Gemäß Ergebnissen bei Tests des Quark-Puzzles ist ein \glqq ansatzloses\grqq{} Puzzlen wenig nachhaltig. Daher ist das Puzzle für Einstiege eher ungeeignet bzw. zu komplex.  

Es hat sich bewährt, dass die Kenntnis der Quarks des Standardmodells der Teilchenphysik und die Forderung nach Farbneutralität von Hadronen einen ausreichenden Ansatz bilden. Der Quark-Inhalt von Nukleonen kann ebenfalls eine zweckmäßige Vorkenntnis darstellen. Im Anschluss kann in einer Erarbeitungsphase das Quark-Puzzle ausgegeben werden (entdeckendes Lernen). Eine mögliche Aufgabenstellung könnte lauten, dass Schüler*innen eigenständig \emph{Teilchen} bauen, den Quarkinhalt der Teilchen festhalten und während des Prozesses nach Eigenschaften der gebauten Teilchen, Gemeinsamkeiten, Unterschiede und weiteren Auffälligkeiten suchen und diese (aus didaktischer Sicht) bestmöglich in Form von Fragen oder Hypothesen verschriftlichen. Die Sammlung von Namen und Auffälligkeiten kann Cloud-basiert oder auf Zetteln erfolgen.  Wir haben außerdem die Erfahrung gemacht, dass Schüler*innen den Namen der gebauten Hadronen erfahren möchten, dies kann mit dem \emph{Hadron-Namer} des CERNs [1] erfolgen. Wenn anstelle der Quarkinhalte der Hadronen deren Namen gesammelt werden, könnten die Multiplett-Darstellungen der Hadronen rekonstruiert werden (z.B. zur Differenzierung). Es bietet sich an, die gesammelten Quarkinhalte/Namen und die herausgearbeiteten Auffälligkeiten im Anschluss, bspw. in einer eher lehrperson-zentrierten Phase zu reflektieren und sich dabei an den aufgestellten Hypothesen zu orientieren. Hierbei können Eigenschaften von Hadronen näher erläutert, konkretisiert und systematisiert werden. Bei uns zeigte sich bei diesem Vorgehen in hoher emotionaler Bezug zu den Inhalten des Vortrags, reges Interesse und kritische Rückfragen. Es konnten außerdem Inhalte vermittelt/veranschaulicht werden, die sonst konzeptuell herausfordernd wären, z.B. Confinement. 

Wird das Quark-Puzzle als Lerngegenstand verwendet, so bilden die Zielgruppe Jugendliche ab der Mittelstufe. Das Quark-Puzzle kann außerdem ohne physikalischen Kontext als \emph{Farbenpuzzle} mit dem Ziel des Baus von \emph{farbneutralen Objekten} dienen. Dies wurde bereits an Fünftklässler*innen erfolgreich getestet. 

\subsection*{Man beachte...}
\begin{itemize}[label=...]
    \item dass auf die Bedeutung der Farbe als Analogie der zugrundeliegenden Symmetrie hingewiesen wird.
    \item dass Quarks als punktförmige Teilchen angesehen werden.
    \item dass es Seequarks gibt, die nicht veranschaulicht werden.
    \item dass es exotische Hadronen mit mehr als drei Valenzquarks gibt.
    \item dass der Isospin unzureichend modelliert werden kann (Drehimpulsalgebra).
\end{itemize}
\subsection*{Kontakt - LHCb Bonn}

\begin{table}[h]
    \centering
    \begin{tabular}{ccc}
        Lukas Julian Exner & Klaas Padeken & Sebastian Neubert  \\
      \textsf{\ding{41} \footnotesize  exner@hiskp.uni-bonn.de} & \textsf{\ding{41} \footnotesize  padeken@hiskp.uni-bonn.de} & \textsf{\ding{41} \footnotesize  neubert@hiskp.uni-bonn.de}
    \end{tabular}
    \label{tab:my_label}
\end{table}

\subsection*{Referenzen}
[1] Adam Morris (2022). \emph{Hadron name builder}. LHCb/CERN. Verfügbar unter:\\ \url{https://hadron-names.web.cern.ch/}
\subsection*{Bilder}
\begin{center}
  \includegraphics[width=0.8\textwidth]{Pictures Quark Puzzle/BA_Übersicht alle Quarks.png}\\ \, \\ 
\includegraphics[width=0.8\textwidth]{Pictures Quark Puzzle/IMG_1824 - Kopie (2).jpg}  
\end{center}

   
\end{document}
